\section{Appendix}
\label{sec:appendix}

\subsection{Analysis of All production traces}
\label{ap:trace}

\begin{figure}[t]
  \centering
  \begin{subfigure}[t]{0.48\linewidth}
    \includegraphics[width=\linewidth]{figs/corr/plot_switch-fixed.pdf}
    \vspace{-5mm}
    \caption{Bursty-group shifts.}
    \label{fig:bursty_shifts}
  \end{subfigure}
  \begin{subfigure}[t]{0.48\linewidth}
    \includegraphics[width=\linewidth]{figs/corr/plot_daytime_correlation-fixed.pdf}
    \vspace{-5mm}
    \caption{Request rate correlation.}
    \label{fig:corr1}
  \end{subfigure}
  \vspace{-2mm}
  \caption{Request patten shifts.}
  \vspace{1mm}
  \label{fig:corr}
\end{figure}

\MyPara{Overview.}
To understand the workload of multi-LLM serving, we analyze four production traces in detail. These traces are collected from representative service providers as summarized in Table~\ref{tab:trace_summary}.
The first two traces are from Hyperbolic~\cite{hyperbolic} and Novita AI~\cite{novita-ai}, two popular LLM inference service providers. They offer inference APIs for a variety of foundation models and also support user-deployed, fine-tuned models. 
The last two traces are from Chatbot Arena~\cite{chatbotarena@arxiv24}, a widely used open-source platform for LLM evaluation. It compares model responses via human preference voting (Arena-Battle) and also provides interfaces for real-time conversations with various models (Arena-Chat).

\MyPara{Bursty-group behaviors.}
The bursty groups in the production traces shift rapidly over time, and this behavior can be quantified through model switches. As shown in Figure~\ref{fig:bursty_shifts}, we plot the number of model switches per hour for each trace. To compute this metric, we treat a model as active if it has received at least one request in the past two minutes. A model switch is counted whenever the set of active models changes. Across the four traces, Novita exhibits the fewest switches, yet still averages 54 switches per hour, meaning the active set changes almost once per minute. In contrast, Arena-Chat, which contains the largest and most diverse collection of models, experiences an average of 766 switches per hour, reflecting that its active set changes every few seconds.

\MyPara{Unpredictability.}
The request rate at a given time is also highly unpredictable in our traces. Figure~\ref{fig:corr1} quantifies predictability by measuring, for each model, the Pearson correlation between its request-rate time series on consecutive days. Across all traces, the correlations cluster near zero, showing that a model’s traffic at a given time of day provides virtually no information about its traffic at the same time on the next day.

\MyPara{Volatile requests.}
Figure~\ref{fig:hour-idle-time} further shows that models frequently alternate between activity and idleness: many models in the Hyperbolic and Novita traces experience more than 40–100 idle intervals per hour (>10 s each), making static reservations wasteful. Complementing this, Figure~\ref{fig:cv-cdf} reports the coefficient of variation (CV), calculated as the standard deviation divided by the mean ($\sigma/\mu$), of requests per minute for each model. Both Hyperbolic and Novita traces contain many models with CV > 1. The Chatbot Arena traces have lower request rates, resulting in smaller CVs overall, but many models still exhibit CVs greater than 0.5. Together, the high frequency of idle intervals and the substantial variability captured by CV indicate that all traces exhibit strong, persistent volatility.

\begin{figure}[t]
  \centering
  \begin{subfigure}[t]{0.48\linewidth}
    \includegraphics[width=\linewidth]{figs/hourly_idle_count_cdf_vertical-fixed.pdf}
    \vspace{-5mm}
    \caption{Number of idle intervals per hour.}
    \label{fig:hour-idle-time}
  \end{subfigure}
  \begin{subfigure}[t]{0.48\linewidth}
    \includegraphics[width=\linewidth]{figs/cv_cdf_all_traces_vertical-fixed.pdf}
    \vspace{-5mm}
    \caption{CV of request rate per minute.}
    \label{fig:cv-cdf}
  \end{subfigure}
  \vspace{-2mm}
  \caption{Volatile requests.}
  \vspace{1mm}
  \label{fig:volatile}
\end{figure}

\subsection{Analysis of Algorithm~\ref{alg:global}}
\label{ap:alg-gloabl}

\subsubsection{KVPR Bound Analysis}
\label{ap:kvpr-bound}

The global model placement algorithm ensures that the maximum KV pressure ratio (KVPR) across all GPUs is bounded by the maximum KVPR in the optimal placement. We give the following analysis.

Let $KVPR_{OPT}$ be the minimum possible maximum KVPR achievable by any optimal placement. Let $KVPR_{max}$ be the maximum KVPR produced by Algorithm \ref{alg:global}. We want to show $KVPR_{max}$ is bounded by $KVPR_{OPT}$.

\MyPara{Bottleneck Analysis:} Focus on the GPU, denoted as $g_{max}$, that achieves the highest KVPR ($KVPR_{max}$) given by Algorithm \ref{alg:global}'s placement. 
Let $m_k$ be the \emph{last} model assigned to this GPU $g_{max}$. Let $W_{before}$ and $S_{before}$ represent the total SLO-weighted request rate and shared KV memory on $g_{max}$ \emph{just before} model $m_k$ was assigned. The final state on this GPU is $KVPR_{max} = (W_{before} +d_k) / (S_{before} - w_k)$, where $d_k$ is the SLO-weighted request rate ($r_k/s_k$) and $w_k$ is the memory weight of model $m_k$.
Similar to Graham~\cite{graham1969bounds}, this proof aims to demonstrate that both the state \emph{before} $m_k$ was added and the contribution \emph{of} $m_k$ are bounded relative to $KVPR_{OPT}$. Specifically, it seeks to establish two conceptual bounds:
\begin{itemize}
    \item \textbf{Bound 1 (Related to state before $m_k$):} The KVPR on $g_{max}$ just before $m_k$'s assignment, $\frac{W_{before}}{S_{before}}$, was the minimum among all GPUs at that moment due to the algorithm's greedy choice. This minimum KVPR is typically related to the average ``pressure'' across the system, which, in turn, is argued to be no larger than the optimal maximum pressure. This suggests the inequality: $W_{before} / S_{before} \le KVPR_{OPT}$
    \item \textbf{Bound 2 (Related to model $m_k$):} The ``pressure'' exerted by the critical model $m_k$ must be handled by the optimal solution. A fundamental lower bound on the optimal solution is the maximum pressure any single model would exert if placed alone on an otherwise empty GPU, \ie, $KVPR_{OPT} \ge d_k/(C-w_k)$.
\end{itemize}
The final step involves integrating these insights to bound $KVPR_{max} = \frac{W_{before} + d_k}{S_{before} - w_k}$. 
Following Graham's proof~\cite{graham1969bounds}, we substitute these into the numerator and get $KVPR_{max}\le KVPR_{OPT}\cdot (1+\frac{C}{S_{g_{max}}-w_k})$.

\subsubsection{TP Support}
\label{ap:tp}
The model placement algorithm in Algorithm~\ref{alg:global} seamlessly integrates models utilizing Tensor Parallelism (TP). We conceptualize a TP model requiring $tp\_size$ GPUs as being composed of $tp\_size$ distinct parts.
For scheduling purposes, we create $tp\_size$ entries in the sorted model list for such a model, assigning each entry $\frac{1}{tp\_size}$ of the original weight and request rate. 
A beneficial property emerges from this decomposition: since these entries have identical $\frac{r_k}{s_k}$ values, they remain adjacent after sorting.
This adjacency increases the likelihood that, as the algorithm iterates, these parts are initially assigned to different GPUs due to rising KVPRs.
To ensure the distribution, if assigning a TP part to the GPU with the minimum KVPR would result in collocating it with another part of the same original model, we instead assign it to the GPU exhibiting the second-lowest KVPR. 
Through this decomposition strategy and modified assignment rule, our algorithm effectively considers and manages the placement of TP models alongside single-GPU models.

\subsection{Elastic Memory Overhead}
\label{sec:elastic-overhead}

\begin{figure}[t]
  \centering    \includegraphics[width=\linewidth]{figs/overhead/mean_latency_comparison-fixed.pdf}
    \vspace{-8mm}
    \caption{Mean latency comparison.
    }
    \label{fig:mean-latency-comparison}
    \vspace{-2mm}
\end{figure}

To stress-test the worst case, we evaluate the system under a constant request rate without idle periods, eliminating opportunities for memory coordination through either time- or space-sharing. In this setting, static partitioning provides the strongest baseline, as it divides memory according to the steady request rate. Figure~\ref{fig:mean-latency-comparison} presents a comparison between our system and static partitioning when serving two Llama-3.2-3B models on a single A100-40G GPU. In this experiment, the global scheduler does not alter placements and the GPU-local scheduler does not reorder requests, since all requests have equal priority. The only additional cost arises from the elastic memory manager, which dynamically maps and unmaps pages, unlike static partitioning that relies on pre-allocation. With optimizations such as buffer preallocation and contiguous layouts, this overhead remains modest: mean TTFT and TPOT increase by only 3–5\% as request rates scale, demonstrating the efficiency of our system even under worst-case conditions.

\subsection{Sensitivity Analysis}
\label{sec:sensitivity}

\begin{figure}[t]
  \centering    \includegraphics[width=\linewidth]{figs/sensitivity_mean_ttft-fixed.pdf}
    \vspace{-8mm}
    \caption{Sensitivity to hyperparameters.}
    \label{fig:sensitivity}
    \vspace{-2mm}
\end{figure}

We analyze the sensitivity of \tool to two key hyperparameters: the model eviction idle threshold and the load monitoring window size with Hyperbolic and Chatbot Arena traces. The results are shown in Figure~\ref{fig:sensitivity}.

\noindent \textbf{Idle Eviction Threshold.} 
Figure~\ref{fig:sensitivity}(a) reports the mean TTFT as we vary the idle threshold, which dictates how long a model must remain inactive before being evicted.
The results reveal a clear convex trade-off. 
When the threshold is too short (\eg, $<40$\,s), the system becomes overly aggressive, evicting models during short inter-arrival gaps. This leads to instability and thrashing, where the system must frequently pay the high penalty of model reactivation for subsequent requests, thereby inflating TTFT. 
Conversely, when the threshold is too long (\eg, $>80$s), idle models hoard GPU memory that could otherwise be harvested. This resource locking prevents the scheduler from placing new active models or expanding the KV cache for concurrent requests, resulting in resource starvation and increased queuing delays. 
Empirically, a threshold of approximately 45 seconds achieves the optimal balance between minimizing reactivation overhead and maximizing resource reclamation.

\noindent \textbf{Load Monitoring Window Size.}
Figure~\ref{fig:sensitivity}(b) examines the impact of the sliding window size used to calculate the moving average of token rates for the KVPR.
This parameter controls the sensitivity of the placement algorithm: a small window makes the scheduler sensitive to transient bursts, while a large window focuses on long-term trends.
The results demonstrate that \tool is generally robust to variations in window size. 
We observe that a window size of approximately 60 seconds provides a stable estimation of memory pressure, effectively smoothing out short-term noise while remaining responsive enough to shift resources during sustained workload changes.

